\documentclass[12pt,a4paper]{article}

% Кодировка и русский язык
\usepackage[utf8]{inputenc}
\usepackage[T2A]{fontenc}
\usepackage[russian,english]{babel}

% Страница и текст
\usepackage[
    a4paper,
    left=3cm,
    right=1.5cm,
    top=2cm,
    bottom=2cm
]{geometry}
\usepackage{setspace}
\usepackage{indentfirst}
\usepackage{enumitem}
\usepackage{titlesec}
\usepackage{microtype}

\onehalfspacing
\setlength{\parindent}{1.25cm}
\setlength{\parskip}{0pt}
\setlist[itemize]{leftmargin=1.25cm}
\setlist[enumerate]{leftmargin=1.25cm}

\titleformat{\section}
    {\bfseries\large}{\thesection.}{0.5em}{}
\titleformat{\subsection}
    {\bfseries\normalsize}{\thesubsection.}{0.5em}{}

% Таблицы и математика
\usepackage{booktabs}
\usepackage{tabularx}
\usepackage{array}
\usepackage{amsmath}
\usepackage{amssymb}

% Ссылки
\usepackage{xcolor}
\usepackage[colorlinks=true,allcolors=blue]{hyperref}

\emergencystretch=2em

\begin{document}

\begin{titlepage}
\begin{center}
\large
Санкт-Петербургский государственный университет
\end{center}

\vspace{3cm}

\begin{center}
\textbf{\large ОТЧЕТ}

\vspace{0.5cm}

по курсовой работе
\end{center}

\vspace{1.5cm}

\begin{center}
\textbf{\LARGE
Подготовка русскоязычного датасета и исследование моделей для чтения по губам
}
\end{center}

\vfill

\begin{flushright}
\begin{minipage}{0.45\textwidth}
\raggedright
Выполнила: студент\\
Тарасова Дарья

\vspace{0.5cm}

Научный руководитель:\\
доцент кафедры ИАС,\\
к. ф.-м. н. Графеева Н. Г.
\end{minipage}
\end{flushright}

\vfill

\begin{center}
Санкт-Петербург\\
2026
\end{center}
\end{titlepage}

\section*{Аннотация}
\addcontentsline{toc}{section}{Аннотация}

В работе продолжен проект по чтению отдельных русских слов по
видеопоследовательностям губ. В третьем семестре был подготовлен
демонстрационный датасет. В четвертом семестре разработан воспроизводимый
конвейер расширения и контроля данных, подготовлены speaker-based split,
mouth ROI и landmark-based preprocessing, а также проведено сравнение
PyTorch-моделей с переносом знаний из lip-reading-модели, обученной на
англоязычном датасете LRW.

Лучшим validation-кандидатом стала frozen LRW-pretrained visual frontend
вместе с pretrained MS-TCN temporal backend. На cross-speaker validation
средний macro-F1 составил $0.5832$. В отложенном post-hoc benchmark на
508 клипах трех новых спикеров получены accuracy
$0.6280 \pm 0.0020$, macro-F1 $0.6124 \pm 0.0040$ и balanced accuracy
$0.6074 \pm 0.0085$. Test ранее просматривался в BiGRU-эксперименте,
поэтому эти числа фиксируются как post-hoc benchmark и не используются
для дальнейшей настройки.

\tableofcontents
\newpage

\section{Введение}

Чтение по губам является задачей компьютерного зрения, в которой по
последовательности кадров видео определяется произнесенная речь или отдельные
слова. Такие методы могут применяться в мультимодальных системах распознавания
речи, при анализе видеоданных и создании вспомогательных технологий.

В данной работе продолжен проект по чтению отдельных русских слов по
видеопоследовательностям области рта. В предыдущем семестре был подготовлен
демонстрационный датасет. В текущем семестре разработан воспроизводимый
конвейер расширения и контроля данных, подготовлены speaker-based split,
mouth ROI и landmark-based preprocessing, а также проведено сравнение
моделей PyTorch с переносом знаний из lip-reading-модели, обученной на
англоязычном датасете LRW.

Главным содержательным результатом стала не одна отдельно взятая модель,
а последовательное исследование, в котором проверялись корректность
обучающего конвейера, влияние дисбаланса, временной структуры, спикерного
разделения, preprocessing и предварительно обученных компонентов.

\section{Постановка задачи}

Цель работы --- исследовать, насколько компактные модели компьютерного
зрения способны классифицировать изолированные русские слова по видео
области рта.

В рамках текущего этапа были поставлены следующие задачи:

\begin{enumerate}
\item расширить датасет по сравнению с демонстрационной версией;
\item автоматизировать повторяющиеся этапы подготовки данных;
\item подготовить единый формат входных данных для нейронной сети;
\item исключить утечку информации между спикерами при разделении данных;
\item реализовать baseline-модель на PyTorch;
\item сравнить temporal-модели и перенос предварительно обученных lip-reading
      признаков;
\item оценить качество по метрикам, устойчивым к дисбалансу классов;
\item зафиксировать ограничения и отрицательные результаты экспериментов.
\end{enumerate}

Задача ограничена десятью словами. Таким образом, рассматриваемая система
не является системой распознавания произвольной русской речи, а представляет
собой контролируемый word-level benchmark.

\section*{История и порядок выполнения работы}
\addcontentsline{toc}{section}{История и порядок выполнения работы}

Работа развивалась поэтапно. Сначала был воспроизведен и проверен простой
baseline, затем отдельно проверялась корректность данных и обучения, после
чего проводилось сравнение temporal-моделей и pretrained-компонентов.
Поэтому notebooks в репозитории не являются независимыми готовыми решениями:
часть из них фиксирует промежуточные эксперименты, отрицательные результаты
или варианты одной и той же модели.

\subsection*{Этап 1. Исходная точка и структура проекта}

В третьем семестре был подготовлен demo-датасет и базовый preprocessing
pipeline. В начале четвертого семестра был создан отдельный проект для
продолжения работы: код preprocessing, локальная структура данных, notebooks
и документация были отделены от материалов третьего семестра. На этом этапе
модель еще не являлась результатом исследования; главной задачей было получить
воспроизводимую основу для дальнейших экспериментов.

\subsection*{Этап 2. Первый PyTorch baseline}

Был добавлен чистый Colab notebook с первым PyTorch baseline. Он включал
восстановление архива, загрузку CSV, Dataset, DataLoader, загрузку 24 кадров,
модель Simple3DCNN, training loop, validation и графики.

Затем были проверены варианты weighted CrossEntropyLoss и
WeightedRandomSampler. Эти эксперименты не дали устойчивого улучшения.
Простая 3D-CNN в основном предсказывала наиболее частый класс. Это показало,
что одной архитектуры и стандартной accuracy недостаточно: нужно учитывать
дисбаланс, temporal structure и speaker-independent evaluation.

\subsection*{Этап 3. Переход к временной модели и аудит ошибок}

Был добавлен controlled FrameCNN+BiGRU baseline. Вместо обработки всего
клипа одной 3D-сверткой модель извлекала признаки из отдельных кадров,
формировала последовательность и передавала ее в двунаправленную GRU.
Параллельно был добавлен model-based clip audit: клипы с большим loss и
неправильными предсказаниями сохранялись в CSV и contact sheets для
визуальной проверки.

Затем были добавлены эксперименты с ResNet18. Сначала ResNet18 обучалась
с нуля, затем использовался frozen pretrained encoder. Отдельно проверялось,
помогает ли balanced training. Эти запуски показали, что на небольшом
датасете обучение большого visual encoder с нуля нестабильно, поэтому
более перспективным направлением стал перенос готовых визуальных признаков.

\subsection*{Этап 4. Проверка самого training pipeline}

В основной notebook был добавлен Tiny Overfit Diagnostic, а затем создан
отдельный notebook для его запуска. Целью было проверить, способен ли код
запомнить небольшой набор, а не получить высокое качество на новых спикерах.

Первый запуск был нестабильным: 16 клипов дали 87.5\%, а 32 клипа ---
65.6\%. После стабилизации эксперимента были зафиксированы seed,
\texttt{num\_workers=0}, кэширование видео, GroupNorm вместо BatchNorm,
отключены dropout, augmentation и weight decay. Были добавлены проверки
путей, labels, конфликтующих дубликатов, движения между кадрами и дисперсии
пикселей.

В стабильном запуске модель достигла 100\% memorization accuracy на 16
клипах на 43-й эпохе и на 32 клипах на 129-й эпохе. Оба диагностических
порога были пройдены. После этого дальнейшее плохое качество full-dataset
baseline можно было рассматривать как проблему обобщения, дисбаланса,
объема и качества данных, а не как очевидную ошибку загрузки или
backpropagation.

\subsection*{Этап 5. Контролируемое сравнение моделей}

После Tiny Overfit были вынесены отдельные controlled notebooks:
FrameCNN+BiGRU на полном speaker-based split, frozen ResNet18 + BiGRU и
balanced frozen-feature experiment. Для этих запусков фиксировались seed,
протокол обучения, validation split и метрики. Сравнение проводилось не
только по accuracy, но и по macro-F1, balanced accuracy, распределению
предсказаний и per-class результатам.

\subsection*{Этап 6. Расширение датасета}

На этапе расширения были добавлены новые кандидаты спикеров. Затем исходные
ссылки и metadata были вынесены из публичной части проекта, чтобы в GitHub
оставались только код, отчеты и необходимые небольшие metadata-файлы.

После этого был подготовлен и оценен Dataset v2, а затем добавлено
multi-seed сравнение. Это позволило проверить, является ли улучшение модели
устойчивым к случайной инициализации, а не результатом одного удачного
запуска.

\subsection*{Этап 7. Pretrained LRW frontend и temporal backends}

Последовательно были добавлены TCN backend, pretrained LRW visual frontend
и финальная тестовая процедура. Затем был проведен отдельный
speaker-independent final benchmark.

Смысл этого этапа состоял в переносе visual frontend, уже обученного на
задаче lip reading, вместо обучения большого visual encoder с нуля. После
этого сравнивались BiGRU и MS-TCN как temporal backend. MS-TCN стал лучшим
кандидатом на cross-speaker validation.

\subsection*{Этап 8. Landmark preprocessing и Dataset v3}

Следующим этапом был добавлен landmark-aligned mouth crop. В отличие от
простого crop по прямоугольнику лица, он использует точки губ и глаз,
сглаживает их во времени и компенсирует наклон головы. Затем были
подготовлены Dataset v3, speaker/sampler ablations и варианты fine-tuning
последнего блока frontend.

На этом этапе сравнивались не только модели, но и варианты входных данных:
обычные mouth crops, landmark crops, horizontal flip и различные составы
спикеров. Это позволило отделить влияние архитектуры от влияния
preprocessing.

\subsection*{Этап 9. Финальное сравнение и контроль риска переобучения}

Затем были проверены перенос pretrained MS-TCN, DC-TCN и temporal masking.
DC-TCN и temporal masking не дали достаточного преимущества и не были
выбраны финальным направлением.

После сравнения моделей работа перешла в фазу аудита и фиксации решения:
был составлен план улучшения модели, проверены границы padded-клипов и их
перекрытия, оценены bounded/moderate timing variants и подготовлен
воспроизводимый post-hoc benchmark.

Финальный post-hoc тест был выполнен после фиксации архитектуры и протокола.
Тестовые клипы ранее просматривались в одном из промежуточных экспериментов,
поэтому результат явно обозначен как post-hoc benchmark, а не как полностью
неиспользованный test set.

\subsection*{Этап 10. Фиксация результатов и документация}

На финальном этапе были сохранены результаты MS-TCN, удалены внутренние
материалы студенческого guide из публичного репозитория и добавлен черновик
основного отчета.

Итоговая последовательность исследования:

\begin{verbatim}
старый preprocessing
 -> первый Simple3DCNN baseline
 -> проверка дисбаланса
 -> FrameCNN+BiGRU
 -> диагностика Tiny Overfit
 -> ResNet и frozen encoder
 -> расширение Dataset v2/v3
 -> LRW-pretrained frontend
 -> BiGRU/TCN/MS-TCN comparison
 -> landmark preprocessing
 -> overlap/data audit
 -> финальный post-hoc benchmark
\end{verbatim}

\section{Результат 3 семестра}

В предыдущем семестре был собран демонстрационный набор данных из 182 клипов,
8 слов и 5 спикеров. Основная цепочка подготовки данных выглядела следующим
образом:

\begin{center}
\begin{verbatim}
video + transcript
 -> audio.wav + cleaned text
 -> WebMAUS forced alignment
 -> TextGrid
 -> words_frames.txt
 -> word-level clips
 -> labels.csv + vocab.txt
 -> validation
\end{verbatim}
\end{center}

В третьем семестре модель распознавания не обучалась. Результатом этапа были
видеоклипы отдельных слов, метаданные и автоматическая проверка
согласованности подготовленного архива.

\section{Данные 4 семестра}

\subsection{Версии датасета}

В работе сохранены несколько версий датасета, поскольку каждая из них
соответствует отдельному этапу экспериментов.

\begin{table}[htbp]
\centering
\begin{tabularx}{\textwidth}{>{\raggedright\arraybackslash}p{0.18\textwidth} >{\raggedright\arraybackslash}X >{\raggedright\arraybackslash}p{0.29\textwidth}}
\toprule
Версия & Назначение & Размер \\
\midrule
Demo 3 семестра & демонстрация конвейера & 182 клипа, 8 слов, 5 спикеров \\
Dataset v1 & первый speaker-based baseline & 955 клипов, 10 слов, 6 спикеров \\
Dataset v2 & добавление двух train-спикеров & 1355 клипов, 10 слов, 8 спикеров \\
Dataset v3 & расширенный train и landmark crops & 1901 train, 105 validation, 508 final-test-кандидатов \\
\bottomrule
\end{tabularx}
\end{table}

Полный набор видеоклипов хранится локально и на Google Drive. В репозитории
GitHub находятся только код, отчеты и небольшие metadata-файлы.

\subsection{Разделение по спикерам}

Для основной оценки моделей используется speaker-based protocol:

\begin{itemize}
\item train --- 1901 клип, 9 спикеров;
\item validation --- 105 клипов, спикер \texttt{spk06};
\item final test --- 508 клипов, спикеры \texttt{spk13}, \texttt{spk14} и \texttt{spk15};
\item словарь --- 10 русских слов.
\end{itemize}

Такое разделение позволяет проверять способность модели переносить
полученные признаки на людей, которых она не видела во время обучения.
Случайное разделение данных используется только в качестве
диагностического контроля.

\subsection{Ограничения данных}

Классы не идеально сбалансированы. Например, в final test слово «есть»
имеет 124 примера, а слово «значит» --- 12. Поэтому одной метрики accuracy
недостаточно для оценки качества.

Кроме того, часть padded-интервалов может пересекаться с соседними словами.
Это проверялось отдельным overlap-аудитом. Результаты моделей не следует
трактовать как доказательство идеальной временной разметки всех клипов.

\section{Pipeline подготовки данных}

\subsection{Исходное видео и субтитры}

Дополнительные видео скачиваются на основании локального manifest-файла.
Ссылки, авторы и оригинальные идентификаторы видео намеренно не публикуются
в репозитории GitHub.

Скрипт \texttt{download\_candidate\_videos.py} получает видео и субтитры.
Длинные записи проходят через \texttt{chunk\_raw\_videos\_for\_maus.py},
который делит их на небольшие chunks. Chunk является техническим фрагментом
исходного видео, а не отдельным словом; он нужен, чтобы WebMAUS обрабатывал
ограниченный по длине аудиофайл.

\subsection{Подготовка WebMAUS}

Скрипт \texttt{prepare\_maus\_inputs.py} читает SRT или текстовую
расшифровку, убирает номера субтитров, таймкоды и служебный мусор,
исправляет повторения rolling subtitles, извлекает из видео mono WAV через
ffmpeg и записывает пару \texttt{audio.wav + cleaned\_text.txt}.

\texttt{batch\_prepare\_maus\_inputs.py} запускает этот этап для набора
файлов. Затем \texttt{batch\_submit\_webmaus\_basic.py} вызывает
\texttt{submit\_webmaus\_basic.py}, отправляет подготовленные audio/text
пары в BAS WebMAUS и получает forced alignment в формате TextGrid.
Пакетная обертка повторяет операцию для многих chunks и изолирует ошибки
отдельных файлов.

\subsection{Временная разметка и word-level clips}

\texttt{get\_phonewords\_frames.py} преобразует интервалы слов из TextGrid
в \texttt{words\_frames.txt}, где для каждого кадра указано активное слово.
\texttt{batch\_generate\_phonewords\_frames.py} применяет этот шаг к
множеству TextGrid.

\texttt{cut\_ru\_clips\_from\_words\_frames.py} использует покадровые
метки и вырезает word-level clips. Пакетный вариант
\texttt{batch\_cut\_ru\_clips\_from\_words\_frames.py} запускает ту же
логику и поддерживает разные варианты временного контекста.

Padded clip содержит небольшой запас кадров до и после границы слова. Это
не новый вид разметки, а тот же word-level clip с контекстом, чтобы начало
и конец артикуляции не обрезались слишком жестко.

\subsection{Labels, quality и mouth ROI}

\texttt{build\_labels\_from\_clips.py} формирует \texttt{labels.csv} по
папкам клипов. \texttt{quality\_check\_clips.py} проверяет читаемость
видео, длительность, число кадров и наличие лица, а также создает report
и contact sheets.

\texttt{create\_mouth\_crops.py} делает обычный mouth ROI. Для модельного
сравнения \texttt{create\_landmark\_mouth\_crops.py} использует landmarks
лица, выравнивание по линии глаз и сглаживание геометрии между кадрами.
\texttt{apply\_manual\_mouth\_crop\_filter.py} применяет rejection list
к labels, сохраняя исходные файлы локально для аудита.

\subsection{Split и упаковка}

\texttt{create\_ml\_splits.py} создает random или speaker-based CSV. Для
основной оценки используется speaker-based split.
\texttt{package\_ml\_dataset.py} собирает ZIP только с клипами, на которые
ссылаются CSV, и записывает POSIX-пути, чтобы архив работал в Linux/Colab.

\section{Вход модели}

Из CSV берутся \texttt{clip\_path} и \texttt{word}. OpenCV читает 24 кадра,
каждый кадр переводится в grayscale и приводится к размеру $96 \times 96$.

Один пример имеет форму:

\begin{verbatim}
[channels, frames, height, width] = [1, 24, 96, 96]
\end{verbatim}

В batch форма имеет вид:

\begin{verbatim}
[batch, channels, frames, height, width] = [B, 1, 24, 96, 96]
\end{verbatim}

Кадры не усредняются в одно число. Их порядок сохраняется, потому что
temporal-часть модели должна видеть движение губ.

\section{Модельные варианты}

\subsection{Simple3DCNN}

Первый baseline использует 3D-свертки, ReLU, pooling и линейный classifier.
3D-свертка видит одновременно пространственные координаты и небольшой
отрезок времени. На несбалансированном наборе модель схлопнулась в наиболее
частый класс. Это полезный диагностический baseline, но не финальная модель.

\subsubsection{Попытки улучшить Simple3DCNN}

Были проверены два способа борьбы с дисбалансом классов.

\begin{enumerate}
\item Для weighted \texttt{CrossEntropyLoss} каждому слову задавался вес,
обратно пропорциональный его частоте в train. В одном из запусков лучшая
validation accuracy составила 0.1238, а к пятой эпохе упала до 0.0571,
поэтому устойчивого улучшения не получилось.
\item При использовании \texttt{WeightedRandomSampler} редкие слова
выбирались чаще при формировании train-батчей. Эксперимент также не дал
стабильного преимущества над обычным baseline.
\end{enumerate}

Эти эксперименты показали, что проблема связана не только с частотами
классов. Простая 3D-CNN также недостаточно хорошо моделировала временную
последовательность движения губ, поэтому далее были проверены
\texttt{FrameCNN + BiGRU} и pretrained temporal-модели.

\subsubsection{Результат Simple3DCNN baseline}

Первый запуск проводился на speaker-based validation из 105 клипов и 10
слов. Модель предсказала все 105 клипов одним наиболее частым классом,
«есть». В validation было 32 клипа со словом «есть», поэтому accuracy
составила $32/105=0.3048$.

\begin{table}[htbp]
\centering
\begin{tabular}{lr}
\toprule
Метрика & Результат \\
\midrule
Accuracy & 0.3048 \\
Macro-F1 & 0.0467 \\
Balanced accuracy & 0.1000 \\
Recall класса «есть» & 1.0000 \\
Recall остальных 9 классов & 0.0000 \\
\bottomrule
\end{tabular}
\end{table}

Accuracy 30.48\% не означает, что модель научилась читать по губам.
Она получена за счет дисбаланса классов и стратегии «всегда предсказывать
самое частое слово». Macro-F1 оказался низким, потому что девять из десяти
классов получили нулевой recall и нулевой F1.

\subsubsection{Диагностический эксперимент Tiny Overfit}

Перед дальнейшим сравнением архитектур был проведен sanity check. Модель
\texttt{FrameCNN+BiGRU} в grayscale обучалась на том же маленьком наборе из
16 или 32 клипов, на котором затем измерялась memorization accuracy. Это не
оценка обобщения и не итоговая accuracy модели. Цель --- проверить загрузку
видео, соответствие \texttt{clip\_path -> word -> class\_id}, вычисление
loss, backward pass и обновление параметров optimizer.

В стабильной версии клипы один раз кэшировались в памяти, использовался
\texttt{num\_workers=0}, фиксировался seed, отключались dropout,
augmentation и weight decay, а tiny-модель использовала GroupNorm вместо
нестабильного на маленьком batch BatchNorm. Дополнительно проверялись
дубликаты тензоров с разными labels, нулевое движение и нулевая дисперсия
пикселей.

\begin{table}[htbp]
\centering
\begin{tabular}{crrl}
\toprule
Набор & Требование & Лучшая accuracy & Результат \\
\midrule
16 клипов & не менее 0.95 & 1.00, эпоха 43 & PASS \\
32 клипа & не менее 0.90 & 1.00, эпоха 129 & PASS \\
\bottomrule
\end{tabular}
\end{table}

В обоих случаях label round trip был корректным, конфликтующих дубликатов
тензоров найдено не было, клипов с нулевым движением и нулевой дисперсией
не было. Итоговый статус --- \texttt{PIPELINE PASS}. Это подтверждает,
что базовый training pipeline способен обучаться и запоминать небольшой
набор, но не доказывает хорошее распознавание новых клипов.

\subsection{Frame CNN + BiGRU}

\begin{verbatim}
grayscale frames
 -> CNN для каждого кадра
 -> sequence [B, 24, feature_dim]
 -> bidirectional GRU
 -> Linear
 -> 10 logits
\end{verbatim}

CNN извлекает визуальные признаки, а BiGRU обрабатывает последовательность
кадров в обоих направлениях. На русском датасете обучаются temporal encoder
и classifier.

\subsection{LRW frontend + BiGRU}

Используется visual frontend, предварительно обученный на LRW. Он включает
временную 3D-свертку и ResNet18. Веса можно заморозить и обучать только
BiGRU и classifier. Такой перенос уменьшает число обучаемых параметров.

\subsection{LRW frontend + pretrained MS-TCN}

Текущий лучший вариант имеет следующую структуру:

\begin{verbatim}
[B, 1, 24, 96, 96]
 -> frozen LRW 3D Conv + ResNet18
 -> [B, 24, 512]
 -> frozen pretrained MS-TCN
 -> [B, 768]
 -> Linear(768, 10)
 -> logits
\end{verbatim}

Обучается только линейная голова $768 \mathbin{\to} 10$. MS-TCN моделирует
temporal context на нескольких временных масштабах и использует знания,
полученные на большом lip-reading-датасете.

\section{Target, loss и обучение}

Строковое слово преобразуется в номер класса через \texttt{word\_to\_idx}.
Например, «будет» отображается в 0. Модель возвращает 10 logits, по одному
на каждый класс.

Используется \texttt{CrossEntropyLoss}. Она сравнивает logits с целевым
\texttt{class\_id} и штрафует модель за низкий score правильного слова.
Для первого baseline также проверялся class-weighted вариант.

Один шаг обучения имеет вид:

\begin{verbatim}
optimizer.zero_grad()
outputs = model(batch_x)
loss = criterion(outputs, batch_y)
loss.backward()
optimizer.step()
\end{verbatim}

\texttt{zero\_grad} очищает старые градиенты, forward получает предсказания,
loss считает ошибку, backward вычисляет градиенты, а \texttt{optimizer.step}
обновляет обучаемые параметры. В transfer MS-TCN меняются только параметры
линейной головы.

\section{Экспериментальный протокол}

Для controlled comparisons использовались фиксированные split, словарь,
preprocessing, seeds и метрики. Validation применялась для выбора варианта,
test был отложен до финального benchmark.

Сохраненные результаты включают loss, accuracy, macro-F1, balanced accuracy,
per-class metrics, confusion matrix, prediction distribution и число
предсказываемых классов.

\section{Результаты}

\subsection{Validation}

\begin{table}[htbp]
\centering
\begin{tabularx}{\textwidth}{>{\raggedright\arraybackslash}Xrrr}
\toprule
Модель & Accuracy & Macro-F1 & Balanced accuracy \\
\midrule
Frozen ImageNet ResNet18 & 0.1587 $\pm$ 0.0440 & 0.1416 $\pm$ 0.0458 & 0.1715 $\pm$ 0.0104 \\
Frozen LRW frontend, old crop & 0.2857 $\pm$ 0.0530 & 0.2736 $\pm$ 0.0290 & 0.3157 $\pm$ 0.0039 \\
Frozen LRW frontend, landmark v3 & 0.5460 $\pm$ 0.0110 & 0.5359 $\pm$ 0.0139 & 0.5787 $\pm$ 0.0108 \\
Pretrained MS-TCN, cross-speaker & --- & \textbf{0.5832} & --- \\
\bottomrule
\end{tabularx}
\end{table}

Horizontal flip дал небольшой положительный эффект. Temporal masking ухудшил
macro-F1 до 0.5391, а DC-TCN получил 0.2834 против 0.3900 у BiGRU control.
Partial fine-tuning \texttt{layer4} не дал устойчивого преимущества.

\subsection{Post-hoc final test}

\begin{table}[htbp]
\centering
\begin{tabular}{crrr}
\toprule
Seed & Accuracy & Macro-F1 & Balanced accuracy \\
\midrule
42 & 0.6280 & 0.6082 & 0.5995 \\
123 & 0.6299 & 0.6131 & 0.6163 \\
2026 & 0.6260 & 0.6160 & 0.6063 \\
\textbf{Mean $\pm$ std} & \textbf{0.6280 $\pm$ 0.0020} & \textbf{0.6124 $\pm$ 0.0040} & \textbf{0.6074 $\pm$ 0.0085} \\
\bottomrule
\end{tabular}
\end{table}

Модель предсказывала все десять классов. Majority baseline для test
составляет 0.2441 accuracy, поэтому модель заметно выше простого выбора
самого частого слова. Однако test имеет post-hoc статус и не должен
использоваться для новых решений по архитектуре.

\section{Обсуждение и ограничения}

Главный положительный вывод состоит в том, что pretrained lip-reading
frontend и temporal backend переносятся на русский word-level benchmark
лучше, чем обычный ImageNet frontend и случайно обучаемая temporal-часть.

Качество ограничено небольшим количеством спикеров, дисбалансом слов,
различиями условий съемки и возможными ошибками временных границ. Результаты
нельзя переносить на свободную русскую речь или считать промышленной системой.

Поскольку final test ранее просматривался в BiGRU-эксперименте, MS-TCN test
benchmark честно помечен как post-hoc. Для полностью независимой оценки
будущей модели потребуется новый test с новыми спикерами, не просмотренный
до фиксации протокола.

\section{Воспроизводимость}

Для фиксации результатов и воспроизводимости экспериментов в проекте
сохранены следующие основные файлы:

\begin{itemize}
\item \texttt{scripts/README.md} --- карта скриптов и описание их статуса
      относительно предыдущего семестра;
\item \texttt{reports/FINAL\_MODEL\_RESULTS.md} --- подробные таблицы
      модельных результатов;
\item \texttt{reports/SEMESTER4\_STATUS.md} --- хронология подготовки
      данных и экспериментов;
\item \texttt{reports/P0\_DATA\_AUDIT.md} --- аудит split, source mapping
      и overlap;
\item \texttt{reports/MODEL\_IMPROVEMENT\_PLAN.md} --- план модельных
      экспериментов и decision gates;
\item \texttt{notebooks/lipreading\_lrw\_pretrained\_mstcn\_posthoc\_final\_test\_colab.ipynb}
      --- notebook для финального post-hoc benchmark.
\end{itemize}

Полный архив видеоданных не хранится в репозитории GitHub. Для запуска
используются два ZIP-архива на Google Drive:

\begin{itemize}
\item \texttt{colab\_lipreading\_dataset\_v3\_train\_val.zip};
\item \texttt{colab\_lipreading\_final\_test\_landmark\_v3.zip}.
\end{itemize}

\section{Заключение}

За четвертый семестр проект был расширен от демонстрационной подготовки
данных до исследовательского ML-конвейера. Были подготовлены воспроизводимые
speaker-based split, quality-control и audit-артефакты, baseline на PyTorch,
а также проведено контролируемое сравнение нескольких архитектур и вариантов
переноса предварительно обученных lip-reading-компонентов.

Лучший вариант с pretrained MS-TCN показал macro-F1 0.5832 на
cross-speaker validation. В post-hoc final-test benchmark средний macro-F1
составил $0.6124 \pm 0.0040$, а средняя accuracy ---
$0.6280 \pm 0.0020$.

Для базового завершения текущего этапа дальнейшее увеличение словаря не
требуется. Основными следующими задачами являются проверка оформления
отчета с научным руководителем и подготовка устного объяснения конвейера
подготовки данных, архитектуры модели, функции потерь, используемых метрик
и ограничений эксперимента. Расширение датасета и создание нового полностью
независимого test-набора могут рассматриваться как дальнейшее развитие работы.

\end{document}
